\chapter{Experiments}
\label{ch:experiments}

This chapter presents two experiments on the OpenJ9 JIT compiler. The first asks whether the
fixed stack-frame budget for stack-allocated objects is a real constraint in practice. The
second asks why so few arrays are stack allocated.

\section{Investigating Stack-Frame Utilization}

OpenJ9 limits the stack-frame space available for stack-allocated objects to a fixed threshold
of \textbf{3000 bytes per method}. When the total size of all eligible candidates exceeds this
limit, escape analysis must remove some of them from the candidate list, and those objects are
allocated on the heap instead. The current implementation decides which candidates to remove
using object size alone, starting with the largest objects (Algorithm~\ref{alg:evict}).

\begin{algorithm}[htbp]
  \caption{Candidate eviction in OpenJ9 (size-based)}
  \label{alg:evict}
  \SetKwProg{Proc}{Procedure}{}{}
  \Proc{\textsc{EvictCandidates}(candidates, MAX)}{
    $\mathit{total} \leftarrow$ size of all candidates on the stack\;
    \While{$\mathit{total} > \mathit{MAX}$}{
      $c \leftarrow$ largest contiguous candidate \tcp*{size only}
      \If{no such $c$}{
        $c \leftarrow$ largest non-contiguous candidate\;
      }
      $\mathit{total} \leftarrow \mathit{total} - \mathit{size}(c)$\;
      mark $c$ as heap-allocated\;
    }
  }
\end{algorithm}

This raises two questions:
\begin{itemize}
  \item Is object size alone a good criterion for choosing which objects to remove?
  \item Does the current heuristic lead to suboptimal stack-frame utilization or suboptimal
        performance on real-world workloads?
\end{itemize}

\subsection{Candidate Eviction: An Example}

Listing~\ref{lst:evict} shows a method with three stack-allocation candidates. $O_1$ is used in
a hot loop, while $O_2$ and $O_3$ are allocated only in a block that runs when \texttt{n} is
zero.

\needspace{15\baselineskip}
\begin{lstlisting}[language=Java, caption={A method whose candidates exceed the budget.}, label={lst:evict}]
double work(int n) {
  Vec acc = new Vec();          // O1, 1600 B
  for (int i = 0; i < n; i++) {
    acc.add(i);                 // hot
  }
  if (n == 0) {                 // cold block
    Cfg a = new Cfg();          // O2, 1000 B
    Log b = new Log();          // O3,  800 B
    a.use(b);
  }
  return acc.sum();
}
\end{lstlisting}

Together the candidates need 3400\,B against a budget of 3000\,B, so at least 400\,B must be
evicted. Table~\ref{tab:evict} compares the current rule with a hot-path-aware rule that evicts
the coldest candidate first, and the smallest among candidates that are equally cold.

\begin{table}[htbp]
  \centering
  \begin{tabular}{@{}llll@{}}
    \toprule
    Rule & Evicts & Stays on stack & Stack space used \\
    \midrule
    Largest-first (current) & $O_1$ & $O_2$, $O_3$ & 1800\,B \\
    Hot-path aware          & $O_3$ & $O_1$, $O_2$ & 2600\,B \\
    \bottomrule
  \end{tabular}
  \caption{Effect of two eviction rules on the method in Listing~\ref{lst:evict}.}
  \label{tab:evict}
\end{table}

The size-based rule evicts $O_1$, the one object that is used on the hot path, and keeps two
objects that are almost never allocated. A rule that takes execution frequency into account keeps
the hot object on the stack and also uses more of the budget. Whether this matters in practice
depends on how often real methods actually exceed the budget, which is what our first experiment
measures.

\section{Experiment 1: Stack-Frame Utilization}

\paragraph{Objective.} Find out whether real-world programs ever need more stack space for
stack-allocated objects than the budget allows.

\paragraph{What we measure.} For every compiled method:
\begin{itemize}
  \item how often the stack-allocation candidates exceed the budget ($\mathit{totalSize} >
        \mathit{MAX}$); and
  \item by how much they exceed it, or how much of the budget is left unused (the
        \emph{deficit}, $\mathit{totalSize} - \mathit{MAX}$).
\end{itemize}

\paragraph{How.} Two static debug counters per method, updated once, when that method is
compiled. From the deficit we compute each method's \emph{utilization}: the share of the
3000\,B budget actually taken by its stack-allocated objects.

\paragraph{Workload.} The DaCapo 9.12-MR1-bach benchmark suite~\cite{blackburn2006dacapo}:
large, varied Java applications whose access patterns are close to those of real-world software.

\subsection{Results}

Figure~\ref{fig:util} summarizes the utilization of every compiled method, per benchmark.

\begin{figure}[htbp]
\centering
\begin{tikzpicture}
\begin{axis}[
    width=0.95\textwidth, height=6cm,
    boxplot/draw direction=y,
    ymin=-0.3, ymax=7,
    xtick={1,2,3,4,5,6,7},
    xticklabels={h2, sunflow, eclipse, fop, jython, lusearch, xalan},
    x tick label style={font=\footnotesize}, y tick label style={font=\footnotesize},
    ylabel={\footnotesize Utilization (\% of 3000\,B)},
    ymajorgrids, grid style={gray!25},
    boxplot={box extend=0.5},
    every boxplot/.style={fill=blue!25, draw=blue!60!black, solid, line width=0.8pt},
    extra y ticks={0}, extra y tick style={grid=none},
]
\addplot+[boxplot prepared={lower whisker=0.0, lower quartile=0.0, median=0.0,
          upper quartile=0.0, upper whisker=0.0}, mark=o, solid]
          coordinates {(1,0.8) (1,1.6) (1,2.13) (1,2.4)};
\addplot+[boxplot prepared={lower whisker=0.0, lower quartile=0.0, median=0.3,
          upper quartile=0.4, upper whisker=0.4}, mark=o, solid]
          coordinates {(2,0.8) (2,1.07) (2,6.4)};
\addplot+[boxplot prepared={lower whisker=0.0, lower quartile=0.0, median=0.4,
          upper quartile=0.8, upper whisker=0.8}, mark=o, solid]
          coordinates {(3,4.27)};
\addplot+[boxplot prepared={lower whisker=0.0, lower quartile=0.0, median=0.0,
          upper quartile=2.13, upper whisker=4.26}, mark=o, solid]
          coordinates {};
\addplot+[boxplot prepared={lower whisker=0.0, lower quartile=0.0, median=0.2,
          upper quartile=0.4, upper whisker=0.4}, mark=o, solid]
          coordinates {(5,1.07)};
\addplot+[boxplot prepared={lower whisker=0.0, lower quartile=0.0, median=0.0,
          upper quartile=0.5, upper whisker=1.06}, mark=o, solid]
          coordinates {};
\addplot+[boxplot prepared={lower whisker=0.0, lower quartile=0.0, median=0.0,
          upper quartile=0.4, upper whisker=0.4}, mark=o, solid]
          coordinates {(7,1.6) (7,4.27)};
\end{axis}
\end{tikzpicture}
\caption{Stack-frame utilization per compiled method. Each point is one method: the share of its
         3000\,B budget used by stack-allocated objects. Box: middle 50\% of methods; line:
         median; whiskers: typical range; circles: outlier methods.}
\label{fig:util}
\end{figure}

The budget is far from being a constraint. The median method uses at most 0.4\% of it (12\,B)
in every benchmark, and in h2, fop, lusearch and xalan at least half of the methods use none of
it at all; in h2 the whole box sits at zero, so at least three quarters of its methods use no
stack space. The largest value observed across all benchmarks is 6.4\% (192\,B), in sunflow. In
other words, no compiled method came close to the 3000\,B limit, and the eviction heuristic of
Algorithm~\ref{alg:evict} did not need to run for any method we observed.

\begin{notebox}
The question raised by the eviction example is therefore moot for this suite: a better eviction
rule cannot help when nothing is evicted. The more pressing problem is the opposite one: stack
frames stay \emph{under-utilized}, because few objects are admitted as candidates in the first
place.
\end{notebox}

\section{Array Stack-Allocation Opportunities}

Stack frames stay under-utilized in most methods, yet arrays are admitted to the candidate list
only under strict conditions, and arrays that fail these checks go to the heap even when they
never escape. This leads to two questions:
\begin{itemize}
  \item Why are arrays most often rejected from the candidate list?
  \item Can we be more aggressive with arrays?
\end{itemize}

\subsection{Eligibility of Arrays}

An array allocation is rejected as a stack-allocation candidate for one of two families of
reasons.

\paragraph{Element class unresolved or unknown.}
\begin{itemize}
  \item \emph{Not known at compile time}: the element class is only decided at run time, for
        example when the array is created reflectively.
  \item \emph{Not resolved}: the element class is known but has not been resolved yet, for
        example because the allocation sits on a path the interpreter never executed before the
        method was compiled.
\end{itemize}
Primitive arrays always pass these checks: their element types are known and loaded at VM
start-up.

\paragraph{Array size.}
\begin{itemize}
  \item \emph{Zero length}.
  \item \emph{Not a constant}: the length is only decided at run time.
  \item \emph{Too large}: the length is above a threshold.
\end{itemize}
These checks exist because a stack frame needs a size fixed at compile time, and a single array
must not use up the whole budget. The question, then, is which of these constraints rejects the
most arrays in real programs.

\section{Experiment 2: Why Arrays Are Rejected}

\paragraph{Objective.} Find the reasons arrays do not make it into the candidate list.

\paragraph{What we measure.}
\begin{itemize}
  \item the percentage of candidate arrays that are actually stack allocated; and
  \item the number of array allocations rejected, for each rejection reason.
\end{itemize}

\paragraph{How.} We added an enumeration listing every possible rejection reason, and one static
debug counter per reason, incremented at the point where the array is rejected. Counters are
organised hierarchically under \texttt{ArrayAlloc}; Listing~\ref{lst:h2counters} shows the raw
output for h2. The reasons map to the labels used in the rest of this section as follows:
\texttt{ArraySizeNotConst} is \emph{dynamic size}, \texttt{ClassUnresolved} is \emph{class
unresolved}, \texttt{ElementClassNotLoadaddr} is \emph{element class not known}, and every other
rejection is counted as \emph{unknown}.

\begin{lstlisting}[language={}, numbers=none, basicstyle=\ttfamily\scriptsize, caption={Static debug counters for h2.}, label={lst:h2counters}]
ArrayAlloc                                              |  1195
ArrayAlloc/ArrayCandidateFailed                         |   471 |  39.41%
ArrayAlloc/ArrayCandidateFailed/ArraySizeNotConst       |   413 |  34.56% |  87.69%
ArrayAlloc/ArrayCandidateFailed/ClassUnresolved         |    31 |   2.59% |   6.58%
ArrayAlloc/ArrayCandidateFailed/ElementClassNotLoadaddr |     2 |   0.17% |   0.42%
ArrayAlloc/ArrayCandidateFailed/Unknown                 |    25 |   2.09% |   5.31%
ArrayAlloc/ArrayCandidateOk                             |   718 |  60.08%
ArrayAlloc/StackAllocated                               |     6 |   0.50%
\end{lstlisting}

\paragraph{Workload.} The DaCapo 9.12-MR1-bach benchmark suite, with 20 iterations per
benchmark. Only compiled methods are counted, so the runs must be long enough for the JIT
compiler to warm up.

\subsection{Results}

Table~\ref{tab:arrays} lists the raw counts. Throughout, the number of array allocations
analysed in a benchmark is the sum of the rejected arrays and the candidate arrays, and the
stack-allocated arrays are counted as a subset of the candidates.\footnote{In every benchmark
the top-level \texttt{ArrayAlloc} counter equals \emph{rejected + candidates + stack allocated}
exactly (for h2, $471 + 718 + 6 = 1195$), which suggests that \texttt{StackAllocated} may be
counted separately from \texttt{ArrayCandidateOk}. If so, the stack-allocation rates below
would change only in the second decimal place (for h2, $6/724$ instead of $6/718$).}

\begin{table}[htbp]
  \centering
  \small
  \setlength{\tabcolsep}{5pt}
  \begin{tabular}{@{}lrrrrrrrr@{}}
    \toprule
    & & \multicolumn{2}{c}{Candidates} & \multicolumn{5}{c}{Rejected} \\
    \cmidrule(lr){3-4} \cmidrule(l){5-9}
    Benchmark & Analysed & Total & Stack & Total & Dyn.\ size & Class unres. & Elem.\ class & Unknown \\
    \midrule
    eclipse  & 2{,}689 & 2{,}163 & 6 & 526 & 392 & 85 & 3 & 46 \\
    fop      & 965     & 730     & 4 & 235 & 180 & 39 & 2 & 14 \\
    h2       & 1{,}189 & 718     & 6 & 471 & 413 & 31 & 2 & 25 \\
    jython   & 920     & 691     & 2 & 229 & 164 & 45 & 2 & 18 \\
    lusearch & 762     & 483     & 0 & 279 & 240 & 27 & 2 & 10 \\
    sunflow  & 1{,}269 & 895     & 9 & 374 & 304 & 49 & 3 & 18 \\
    xalan    & 1{,}129 & 733     & 4 & 396 & 282 & 71 & 2 & 41 \\
    \midrule
    All      & 8{,}923 & 6{,}413 & 31 & 2{,}510 & 1{,}975 & 347 & 16 & 172 \\
    \bottomrule
  \end{tabular}
  \caption{Array allocations in compiled methods of the DaCapo 9.12-MR1 benchmarks.}
  \label{tab:arrays}
\end{table}

\paragraph{Candidates are rarely stack allocated.} Most array allocations, \CandRate\% overall,
do pass the candidacy checks. Yet very few of these candidates end up on the stack
(Figure~\ref{fig:stackarrays}): across all benchmarks only \StackCount{} of \CandCount{}
candidate arrays (\StackRateOfCand\%) are stack allocated, and lusearch stack-allocates none.
Passing the candidacy checks is therefore far from sufficient; the arrays are lost at a later
stage of escape analysis.

\begin{figure}[htbp]
  \centering
  \begin{tikzpicture}
    \begin{groupplot}[
        group style={group size=2 by 1, horizontal sep=2.0cm, yticklabels at=edge left},
        benchrows, benchyticks, height=6.2cm, /pgf/bar width=8pt, xmin=0]

      \nextgroupplot[xbar stacked, width=7.4cm, xmax=\OkXmax,
                     title={Candidate arrays (count)},
                     legend columns=2, legend style={at={(1.15,1.18)}, anchor=south}]
      \addplot[seg=cheap] table[x=okheap, y=ypos]{\benchtab};
      \addplot[seg=cstack, endlab] table[x=okstack, y=ypos, meta=oklab]{\benchtab};
      \legend{Not stack allocated (heap), Stack allocated}

      \nextgroupplot[xbar, /pgf/bar shift=0pt, width=6.6cm, xmax=\StackXmax,
                     title={\% of candidates stack allocated}]
      \addplot[fill=cstack, draw=none, endlab,
               restrict expr to domain={\coordindex}{0:\LastBenchIdx}]
        table[x=stackpct, y=ypos, meta=stacklab]{\benchtab};
      \addplot[fill=cstackdark, draw=none, endlab,
               restrict expr to domain={\coordindex}{\AllIdx:\AllIdx}]
        table[x=stackpct, y=ypos, meta=stacklab]{\benchtab};
    \end{groupplot}
  \end{tikzpicture}
  \caption{Stack allocation of candidate arrays. Left: candidate arrays per benchmark, split into
           heap and stack allocated (the stack-allocated part is too small to be visible). Right:
           the share of candidates that is stack allocated.}
  \label{fig:stackarrays}
\end{figure}

\paragraph{Dynamic size dominates the rejections.} Figure~\ref{fig:reasons} breaks the rejected
arrays down by reason. A size that is not a compile-time constant is by far the most common
reason, accounting for \TopReasonShare\% of all rejections (1{,}975 of \FailedCount{}), and for
between 71\% and 88\% in every individual benchmark. An unresolved element class is a distant
second at 13.8\% overall (6.6\% to 19.7\% per benchmark). Arrays whose element class is not
known at compile time are rare (16 in total), consistent with reflective array creation being
uncommon in hot code.

\begin{figure}[htbp]
  \centering
  \begin{tikzpicture}
    \begin{axis}[benchrows, benchyticks, xbar stacked, /pgf/bar width=8pt,
                 width=0.86\textwidth, height=6.4cm, xmin=0, xmax=100,
                 xlabel={\% of rejected array allocations},
                 legend columns=4, legend style={at={(0.5,1.03)}, anchor=south}]
      \addplot[seg=csize,    inseg] table[x=p_size,    y=ypos, meta=p_size]{\benchtab};
      \addplot[seg=cclass,   inseg] table[x=p_class,   y=ypos, meta=p_class]{\benchtab};
      \addplot[seg=celem,    inseg] table[x=p_elem,    y=ypos, meta=p_elem]{\benchtab};
      \addplot[seg=cunknown, inseg] table[x=p_unknown, y=ypos, meta=p_unknown]{\benchtab};
      \addplot[draw=none, fill=none, forget plot, endlab]
        table[x expr=0.001, y=ypos, meta=failedlab]{\benchtab};
      \legend{Dynamic size, Class unresolved, Element class not known, Unknown}
    \end{axis}
  \end{tikzpicture}
  \caption{Reasons why array allocations are rejected as stack-allocation candidates; $n$ is the
           number of rejected arrays in each benchmark.}
  \label{fig:reasons}
\end{figure}

\subsection{Discussion}

The two experiments point in the same direction. The per-method budget is almost never the
limiting factor; what limits stack allocation of arrays is, first, that 28\% of all array allocations
are rejected as candidates, mostly because their size is not a compile-time
constant, and second, that almost none of the arrays that \emph{are} candidates end up on the
stack. Both gaps are targets for the research directions in the next chapter.
